\notechapter{N-20220720}{Counting Principles, Binomial Identities, and Mixed Problem Solving}

\section{Counting principles}

The note begins with the standard counting principles:
\begin{description}
  \item[Addition principle.] If a task can be done in disjoint cases with \(m_1,m_2,\ldots,m_n\) possibilities, then the total number is \(m_1+\cdots+m_n\).
  \item[Multiplication principle.] If a task is done in stages with \(m_1,m_2,\ldots,m_n\) choices at each stage, then the total number is \(m_1m_2\cdots m_n\).
  \item[Inclusion--exclusion.] For sets \(A_1,\ldots,A_n\), one adds single sizes, subtracts pairwise intersections, adds triple intersections, and so on.

\end{description}

A Venn diagram in the note reviews the three-set formula:
\[
\begin{aligned}
|A\cup B\cup C|&=|A|+|B|+|C|\\
&\quad-|A\cap B|-|A\cap C|-|B\cap C|+|A\cap B\cap C|.
\end{aligned}
\]

\section{A divisibility example}

One problem asks how many positive integers not exceeding \(2021\) are neither multiples of \(6\) nor multiples of \(8\).  The excluded set is
\[
  A\cup B,
\]
where \(A\) is the set of multiples of \(6\), and \(B\) is the set of multiples of \(8\).  Since
\[
  \operatorname{lcm}(6,8)=24,
\]
we have
\[
  |A|=336,\qquad |B|=252,\qquad |A\cap B|=84.
\]
Thus the number excluded is
\[
  336+252-84=504,
\]
and the answer is
\[
  2021-504=1517.
\]
This is exactly what inclusion--exclusion is for: avoid double-counting the common multiples.

\section{Permutations and combinations}

The page distinguishes arrangements and selections:
\[
  A_n^m=\frac{n!}{(n-m)!},\qquad
  C_n^m=\frac{n!}{m!(n-m)!}.
\]
A sequence problem asks for the number of \(10\)-term sequences with \(x+y=10\) and \(x-y=4\), giving \(x=7,y=3\).  Therefore the number of sequences is
\[
  \binom{10}{3}=120.
\]
The thought process is: once the number of one kind of symbol is fixed, the sequence is determined by choosing its positions.

\section{Binomial theorem}

The binomial theorem is
\[
  (a+b)^n=\sum_{k=0}^n\binom nk a^{n-k}b^k.
\]
Several standard identities follow immediately:
\[
  \sum_{k=0}^n\binom nk=2^n,
\]
\[
  \sum_{k=0}^n(-1)^k\binom nk=0,
\]
and by differentiating
\[
  (1+x)^n
\]
and setting \(x=1\),
\[
  \sum_{k=0}^n k\binom nk=n2^{n-1}.
\]
The page writes these identities as examples of how a coefficient formula becomes a counting tool.

\section{Polynomial coefficients}

A worked example asks for the constant term or a particular coefficient in an expansion such as
\[
  (2x-2y)^5.
\]
The term involving \(x^{5-k}y^k\) is
\[
  \binom5k(2x)^{5-k}(-2y)^k.
\]
Thus the coefficient of \(x^2y^3\) comes from \(k=3\):
\[
  \binom53 2^2(-2)^3=-320.
\]
The note's point is that coefficient problems become simple once the target exponent determines \(k\).

\section{Further contest examples}

The later pages mix counting with algebra, inequalities, and geometry.  There are examples involving choosing books from different subjects, solving inequalities by sign charts, and geometry problems involving similar triangles.  The unifying skill is case management: one must decide whether the problem is a sum of cases, a product of stages, or a coefficient extraction.


\section{Binomial identities by stories}

The identity
\[
  \binom n k=\binom{n-1}k+\binom{n-1}{k-1}
\]
comes from choosing a committee of size \(k\) from \(n\) people.  Either a distinguished person is not chosen, or that person is chosen.  This kind of proof is often clearer than algebraic manipulation because it explains what the two terms mean.

\section{When to use multiplication and when to use addition}

The addition principle applies to disjoint alternatives.  The multiplication principle applies to successive independent choices.  Many counting mistakes come from using multiplication when cases overlap, or using addition when a task has stages.  A good habit is to write a sentence: ``first choose ..., then choose ...'' for multiplication, and ``either ..., or ...'' for addition.

\section{Expanded synthesis and advanced context}

Elementary counting is a first meeting with combinatorial species and generating functions.  Addition corresponds to disjoint union, multiplication corresponds to Cartesian product, and binomial coefficients are the coefficients of a generating function.  This is why the binomial theorem is not just algebra: it is also a counting machine.


\paragraph{Expanded synthesis.}
The number-theoretic content of this chapter is best understood as arithmetic seen through algebraic structure. The earlier sections (`Counting principles`, `A divisibility example`, `Permutations and combinations`, `Binomial theorem`, `Polynomial coefficients`) compute divisibility, congruences, residues, prime factorizations, or finite-field examples. The deeper language is ideals, quotient rings, unit groups, valuations, and field extensions.

Divisibility in $\mathbb Z$ is containment of principal ideals; congruence modulo $n$ is equality inside the quotient ring $\mathbb Z/n\mathbb Z$; invertibility modulo $n$ means being a unit; the Chinese remainder theorem is a product decomposition of quotient rings. Prime factorization supplies coordinates through valuations:
\[
  n=\prod_p p^{v_p(n)}.
\]
Gcd and lcm are coordinatewise minimum and maximum in these valuation coordinates.

This algebraic reading explains why elementary tricks scale upward. Euclidean algorithms become the theory of Euclidean domains and PIDs; divisor sums become Dirichlet convolution and Euler products; roots of unity become cyclotomic fields and Galois actions; finite polynomial arithmetic becomes finite-field theory. The original computations should therefore be kept as the algorithmic core, while the advanced language identifies the structure that makes the algorithms work.
